\documentclass[a4paper, 12pt]{report}

\usepackage{amsmath}
\usepackage{amssymb}
\usepackage[margin=1in]{geometry}
\usepackage{hyperref}
\usepackage{graphicx}
\usepackage{subcaption}
\usepackage{float}
\usepackage{array}
\usepackage{booktabs}
\setcounter{secnumdepth}{3}

\renewcommand{\thesection}{\arabic{section}}
\renewcommand{\thesubsection}{\arabic{section}.\arabic{subsection}}

\setlength{\parindent}{0pt}
\setlength{\parskip}{6pt}

\begin{document}

\chapter*{Signal Analysis}

Signal analysis is performed in this chapter. The aim is to first provide a brief introduction
to the Discrete Fourier Transform (DFT) in the context of its application to MultiFrequency analysis
for Flutter Derivative Identification. This is then applied to single- and multi-frequency signal extraction.
Finally, potential sources of error and methods to mitigate them are discussed.

This chapter is organised as follows:

\begin{flushleft}
\hspace{2cm}1. Introduction to DFT \\
\hspace{2cm}2. Single Frequency \\
\hspace{2cm}3. Multi-Frequency \\
\end{flushleft}

% ============================================================
\section{Introduction to DFT}
% ============================================================

The Fourier Transform is a mathematical tool used to convert a signal from the time domain to
the frequency domain. Since real-world data consists of discrete numbers rather than continuous
mathematical functions, the transform cannot be applied directly. Therefore, the Discrete
Fourier Transform (DFT) is used in such cases. The DFT converts time-varying data into the
frequency domain, and its inputs are the time-varying data and the sampling frequency. The
sampling frequency defines how many samples are taken per second; in the CFD context used in
this thesis, it is equal to $1/\Delta t$, where $\Delta t$ is the time step.

The Fast Fourier Transform (FFT) is derived from the DFT and offers significantly greater
computational efficiency. Several open-source tools are available to perform FFT computations,
one of which is provided by the \texttt{scipy.fft} module in Python. Since Python is used for
post-processing throughout this thesis, the SciPy FFT tool is employed to convert the obtained
aerodynamic forces and moments into the frequency domain.

Since the DFT operates on a discrete set of values, it is subject to certain inherent
limitations that can affect the quality of the transformation and the usability of its results.
It is therefore necessary to study these issues before applying the FFT to flutter derivative
calculations. This is the focus of the present chapter: to examine the issues associated with
the DFT and to investigate its efficient application in the multi-frequency method for flutter
derivative identification.

The accuracy of extracting the amplitude of a particular signal using the FFT depends on two
parameters:

\begin{enumerate}
    \item The frequency of the signal to be extracted.
    \item The frequency bin of the data.
\end{enumerate}

The frequency bin refers to the discrete set of frequency values produced by the FFT during
domain conversion. If the frequency of the signal to be extracted coincides with one of the
values in the frequency bin, the extracted amplitude will be accurate and free from
discretisation errors.

\vspace{0.3cm}
\textbf{Calculation of Frequency Bin}

% TODO: Include derivation/formula for frequency bin calculation

% TODO: Insert figure of frequency bin

\vspace{0.3cm}

The spacing between two adjacent frequencies in the bin is known as the \textit{bin spacing},
and the \textit{frequency resolution} of an FFT refers to how small this spacing is. A smaller
bin spacing corresponds to a higher frequency resolution. As shown by the formula, the
resolution depends on the length of the data record.

If the signal frequency does not coincide with a frequency bin value, the amplitude obtained
from the FFT will contain an error. This effect is known as \textit{spectral leakage}: since
no bin directly corresponds to the signal frequency, its amplitude is distributed across two
adjacent bins. Therefore, using a longer data record increases the accuracy of the DFT by
reducing the bin spacing and mitigating spectral leakage.

Another important issue to be aware of is \textit{aliasing}. The maximum signal frequency that
can be extracted using an FFT is limited by the Nyquist frequency, given by:

\begin{equation}
    f_{\text{Nyquist}} = \frac{f_s}{2}
\end{equation}

where $f_s$ is the sampling frequency. The implications of this limit are twofold. First, if
the frequency of interest lies above the Nyquist frequency, it cannot be extracted. Second,
even if the frequency of interest lies within the Nyquist limit, any other signal with a
frequency above it will be folded back into the valid frequency range by the FFT. If such a
folded component does not fall on a bin, it will introduce additional spectral leakage. While
the probability of this aliased component coinciding exactly with a frequency of interest is
low, this effect should be kept in mind when using the FFT in practice.

These two effects combined impose restrictions on the frequencies that can be extracted and
the accuracy with which they can be obtained, in various forms:

\begin{enumerate}
    \item As discussed, if the frequency of interest does not have a finite decimal
          representation or does not fall on a frequency bin, spectral leakage will introduce
          error.
    \item The length of the data record determines a frequency bin that may or may not match
          the frequency of interest.
    \item The frequency of interest must lie within the Nyquist frequency.
    \item If random noise such as white noise is present in the data, the extracted amplitude
          will contain error even if the signal frequency falls on a bin. This is because random
          noise excites many frequencies simultaneously, and any overlap with the frequency range
          of the signal of interest will disturb the extracted amplitude. Since the noise is
          random, a large bin spacing further compounds this error through spectral leakage.
\end{enumerate}

\vspace{0.5cm}

% ============================================================
\section{Single Frequency}
% ============================================================

This section studies the effects discussed in Section~1 as applied to a single-frequency
signal. A sample sine wave is generated with discrete points, passed to the FFT tool from the
SciPy module, and the resulting amplitudes are analysed. Since the amplitude of the input
signal is known, it is compared against the FFT output and the error is examined for the
following cases:

\begin{enumerate}
    \item Single Frequency — Bin Match
    \item Single Frequency — Non-Bin Match
    \item Single Frequency — Bin Match with Noise
    \item Frequency Bin Resolution on Non-Matching Bin
\end{enumerate}

The input signal parameters used throughout this section are as follows:

\begin{itemize}
    \item Total time: 10 seconds
    \item Sampling rate: 500 Hz (i.e.\ each second is divided into 500 discrete samples)
\end{itemize}

The resulting frequency bins are:

\textbf{Frequency Bin:} $0 \mid 0.1 \mid 0.2 \mid 0.3 \mid 0.4 \mid 0.5 \mid 0.6 \mid \cdots \mid 250.0$ Hz

Note that the maximum frequency is 250~Hz, which is the Nyquist limit for this configuration.

% ------------------------------------------------------------
\subsection{Single Frequency — Bin Match}
% ------------------------------------------------------------

A signal with a frequency of 0.4~Hz and an amplitude of 1.0 is used as input. Since 0.4~Hz
falls directly on a frequency bin, no spectral leakage is expected.

\begin{figure}[htbp]
    \centering
    \begin{subfigure}[b]{0.45\textwidth}
        \centering
        \includegraphics[width=\textwidth]{Chap_3/Figures/SingleFrequency/1_Input_Signal.png}
        \caption{Input signal at 0.4~Hz.}
        \label{fig:sf_binmatch_input}
    \end{subfigure}
    \hfill
    \begin{subfigure}[b]{0.45\textwidth}
        \centering
        \includegraphics[width=\textwidth]{Chap_3/Figures/SingleFrequency/1_Spectrum.png}
        \caption{FFT amplitude spectrum — bin match.}
        \label{fig:sf_binmatch_spectrum}
    \end{subfigure}
    \caption{Single-frequency bin-match case: input signal and corresponding FFT spectrum.}
    \label{fig:sf_binmatch}
\end{figure}

The amplitude output from the FFT is 1.00, which matches the input signal amplitude exactly.
This confirms that when the signal frequency falls on a frequency bin, no discretisation error
is introduced.

% ------------------------------------------------------------
\subsection{Single Frequency — Non-Bin Match}
% ------------------------------------------------------------

In this case, the input signal frequency is changed to 0.45~Hz, which does not coincide with
any frequency bin. Since 0.45~Hz lies midway between the 0.4~Hz and 0.5~Hz bins, the FFT
assigns its amplitude to the nearest neighbouring bin.

\begin{figure}[htbp]
    \centering
    \begin{subfigure}[b]{0.45\textwidth}
        \centering
        \includegraphics[width=\textwidth]{Chap_3/Figures/SingleFrequency/2_Input_Signal.png}
        \caption{Input signal at 0.45~Hz.}
        \label{fig:sf_nonbinmatch_input}
    \end{subfigure}
    \hfill
    \begin{subfigure}[b]{0.45\textwidth}
        \centering
        \includegraphics[width=\textwidth]{Chap_3/Figures/SingleFrequency/2_Spectrum.png}
        \caption{FFT amplitude spectrum — non-bin match.}
        \label{fig:sf_nonbinmatch_spectrum}
    \end{subfigure}
    \caption{Single-frequency non-bin-match case: input signal and corresponding FFT spectrum showing spectral leakage.}
    \label{fig:sf_nonbinmatch}
\end{figure}

The extracted amplitude is 0.653, corresponding to an error of 34.70\%. This demonstrates
that when the signal frequency does not fall on a frequency bin, spectral leakage causes a
significant reduction in accuracy.

% ------------------------------------------------------------
\subsection{Single Frequency — Bin Match with Noise}
% ------------------------------------------------------------

In this case, the input signal is kept at 0.4~Hz. Additionally, white noise with an amplitude
of 0.2 is superimposed on the signal to assess its effect on the extracted amplitude.

\begin{figure}[htbp]
    \centering
    \begin{subfigure}[b]{0.45\textwidth}
        \centering
        \includegraphics[width=\textwidth]{Chap_3/Figures/SingleFrequency/3_Input_Signal.png}
        \caption{Input signal at 0.4~Hz with white noise.}
        \label{fig:sf_noise_input}
    \end{subfigure}
    \hfill
    \begin{subfigure}[b]{0.45\textwidth}
        \centering
        \includegraphics[width=\textwidth]{Chap_3/Figures/SingleFrequency/3_Spectrum.png}
        \caption{FFT amplitude spectrum — bin match with noise.}
        \label{fig:sf_noise_spectrum}
    \end{subfigure}
    \caption{Single-frequency bin-match case with added white noise: input signal and corresponding FFT spectrum.}
    \label{fig:sf_noise}
\end{figure}

The extracted amplitude is 0.99, corresponding to an error of only 0.04\%. Although noise does
introduce a small error, its effect is considerably less severe than that of spectral leakage,
provided the signal frequency falls on a frequency bin.

% ------------------------------------------------------------
\subsection{Frequency Bin Resolution on Non-Matching Bin}
% ------------------------------------------------------------

% TODO: Add content for this subsection

\vspace{0.5cm}

% ============================================================
\section{Multi-Frequency}
% ============================================================

This section extends the analysis of the previous section to multi-frequency signals. The
primary objectives are to determine whether increasing the number of frequency components
affects accuracy, to examine the impact of a frequency not falling on a bin in a
multi-component signal, and to assess the effect of noise in such cases.

The data length and frequency bins are kept identical to those used in the single-frequency
analysis:

\begin{itemize}
    \item Total time: 10 seconds
    \item Sampling rate: 500 Hz
\end{itemize}

\textbf{Frequency Bin:} $0 \mid 0.1 \mid 0.2 \mid 0.3 \mid 0.4 \mid 0.5 \mid 0.6 \mid \cdots \mid 250.0$ Hz

The following cases are examined:

\begin{enumerate}
    \item Multi-Frequency — All Components Bin-Matched
    \item Multi-Frequency — Bin Match with Noise
    \item Multi-Frequency — One Component Non-Bin-Matched
\end{enumerate}

% ------------------------------------------------------------
\subsection{Multi-Frequency — All Components Bin-Matched}
% ------------------------------------------------------------

In this case, all input frequencies are chosen to fall on frequency bins, and the amplitude
accuracy of the FFT output is verified for each component.

\begin{table}[htbp]
\centering
\caption{FFT amplitude detection results for a multi-frequency bin-matched signal.}
\label{tab:mf_binmatch}
\begin{tabular}{cccc}
\toprule
\textbf{Input Frequency (Hz)} & \textbf{Input Amplitude} & \textbf{FFT Amplitude} & \textbf{Error (\%)} \\
\midrule
0.40  & 1.0 & 1.000 & 0.00 \\
0.80  & 4.0 & 4.000 & 0.00 \\
1.00  & 3.0 & 3.000 & 0.00 \\
1.20  & 1.0 & 1.000 & 0.00 \\
1.40  & 2.0 & 2.000 & 0.00 \\
\midrule
\multicolumn{3}{r}{\textbf{Overall Error}} & \textbf{0.00\%} \\
\bottomrule
\end{tabular}
\end{table}

\begin{figure}[htbp]
    \centering
    \begin{subfigure}[b]{0.45\textwidth}
        \centering
        \includegraphics[width=\textwidth]{Chap_3/Figures/MultiFrequency/1_Input_Signal.png}
        \caption{Multi-frequency input signal — all bin-matched.}
        \label{fig:mf_binmatch_input}
    \end{subfigure}
    \hfill
    \begin{subfigure}[b]{0.45\textwidth}
        \centering
        \includegraphics[width=\textwidth]{Chap_3/Figures/MultiFrequency/1_Spectrum.png}
        \caption{FFT amplitude spectrum — all components bin-matched.}
        \label{fig:mf_binmatch_spectrum}
    \end{subfigure}
    \caption{Multi-frequency bin-matched case: input signal and corresponding FFT spectrum with zero error across all components.}
    \label{fig:mf_binmatch}
\end{figure}

As seen in Table~\ref{tab:mf_binmatch} and Figure~\ref{fig:mf_binmatch}, all five frequency
components are extracted with zero error. This confirms that when all signal frequencies fall
on frequency bins, the FFT produces accurate results regardless of the number of frequency
components present.

% ------------------------------------------------------------
\subsection{Multi-Frequency — Bin Match with Noise}
% ------------------------------------------------------------

The same multi-frequency signal is used here, with the addition of white noise at an amplitude
of 0.2. The aim is to assess whether the noise effect observed in the single-frequency case
persists under multi-frequency conditions.

\begin{table}[htbp]
\centering
\caption{FFT amplitude detection results for a multi-frequency bin-matched signal with added white noise.}
\label{tab:mf_binmatch_noise}
\begin{tabular}{cccc}
\toprule
\textbf{Input Frequency (Hz)} & \textbf{Input Amplitude} & \textbf{FFT Amplitude} & \textbf{Error (\%)} \\
\midrule
0.40  & 1.0 & 1.000 & 0.04 \\
0.80  & 4.0 & 4.000 & 0.15 \\
1.00  & 3.0 & 3.000 & 0.15 \\
1.20  & 1.0 & 1.000 & 0.09 \\
1.40  & 2.0 & 2.000 & 0.07 \\
\midrule
\multicolumn{3}{r}{\textbf{Overall Error}} & \textbf{0.12\%} \\
\bottomrule
\end{tabular}
\end{table}

\begin{figure}[htbp]
    \centering
    \begin{subfigure}[b]{0.45\textwidth}
        \centering
        \includegraphics[width=\textwidth]{Chap_3/Figures/MultiFrequency/2_Input_Signal.png}
        \caption{Multi-frequency input signal with white noise.}
        \label{fig:mf_noise_input}
    \end{subfigure}
    \hfill
    \begin{subfigure}[b]{0.45\textwidth}
        \centering
        \includegraphics[width=\textwidth]{Chap_3/Figures/MultiFrequency/2_Spectrum.png}
        \caption{FFT amplitude spectrum — bin match with noise.}
        \label{fig:mf_noise_spectrum}
    \end{subfigure}
    \caption{Multi-frequency bin-matched case with added white noise: input signal and corresponding FFT spectrum.}
    \label{fig:mf_noise}
\end{figure}

As shown in Table~\ref{tab:mf_binmatch_noise} and Figure~\ref{fig:mf_noise}, the introduction
of white noise produces small but non-zero errors across all frequency components. Consistent
with the single-frequency case, the errors remain small relative to those caused by spectral
leakage, provided all frequencies fall on bins.

% ------------------------------------------------------------
\subsection{Multi-Frequency — One Component Non-Bin-Matched}
% ------------------------------------------------------------

In this case, the effect of a single frequency component not falling on a frequency bin is
investigated within a multi-frequency signal. Specifically, the analysis examines whether the
resulting error is isolated to the non-matching frequency or whether it propagates to
neighbouring frequency components as well.

To investigate this, a multi-frequency signal identical to the previous case is used, with an
additional component at 1.05~Hz introduced. Since the frequency bins are located at 1.0~Hz
and 1.1~Hz, this frequency falls exactly midway between two adjacent bins, making it a
worst-case spectral leakage scenario.

\begin{table}[htbp]
\centering
\caption{FFT amplitude detection results for a multi-frequency signal with one non-bin-matched component at 1.05~Hz.}
\label{tab:mf_nonbinmatch}
\begin{tabular}{cccc}
\toprule
\textbf{Input Frequency (Hz)} & \textbf{Input Amplitude} & \textbf{FFT Amplitude} & \textbf{Error (\%)} \\
\midrule
0.40  & 1.0 & 1.0095 &  0.95 \\
0.80  & 4.0 & 3.8793 &  3.02 \\
1.00  & 3.0 & 2.4999 & 16.67 \\
1.05  & 1.0 & 0.9654 &  3.45 \\
1.10  & 1.0 & 0.2145 & 78.55 \\
1.40  & 2.0 & 2.0152 &  0.76 \\
\midrule
\multicolumn{3}{r}{\textbf{Overall Error}} & \textbf{17.23\%} \\
\bottomrule
\end{tabular}
\end{table}

\begin{figure}[htbp]
    \centering
    \begin{subfigure}[b]{0.45\textwidth}
        \centering
        \includegraphics[width=\textwidth]{Chap_3/Figures/MultiFrequency/3_Input_Signal.png}
        \caption{Multi-frequency input signal with one non-bin-matched component.}
        \label{fig:mf_nonbinmatch_input}
    \end{subfigure}
    \hfill
    \begin{subfigure}[b]{0.45\textwidth}
        \centering
        \includegraphics[width=\textwidth]{Chap_3/Figures/MultiFrequency/3_Spectrum.png}
        \caption{FFT amplitude spectrum showing spectral leakage from the non-bin-matched component.}
        \label{fig:mf_nonbinmatch_spectrum}
    \end{subfigure}
    \caption{Multi-frequency case with one non-bin-matched component at 1.05~Hz: input signal and FFT spectrum illustrating leakage into neighbouring bins.}
    \label{fig:mf_nonbinmatch}
\end{figure}

As shown in Table~\ref{tab:mf_nonbinmatch} and Figure~\ref{fig:mf_nonbinmatch}, spectral
leakage from the 1.05~Hz component does not remain isolated to that frequency alone. It also
corrupts the amplitudes of neighbouring frequency components, particularly 1.00~Hz and
1.10~Hz, which exhibit errors of 16.67\% and 78.55\% respectively.

A key observation is that the magnitude of this leakage-induced error decreases with
increasing spectral distance from the non-bin-matched frequency. This spatial decay of leakage
is also visible in the FFT spectrum, where the spurious disturbance diminishes progressively
at frequencies further from 1.05~Hz. This behaviour underlines the importance of ensuring that
all frequency components in a multi-frequency signal fall on frequency bins before applying the
FFT.

\end{document}